% 原PDF528–533／印刷497–502；全部原页已目视核对。
\clearpage\section*{计算机代数}\phantomsection\label{app:computer-algebra}
\addcontentsline{toc}{section}{计算机代数}
\begingroup
\renewcommand{\thesection}{\arabic{section}}
\setcounter{section}{12}\setcounter{subsection}{6}\setcounter{theorem}{0}
\sourcepageanchor{528}
计算机代数系统（computer algebra system）能够对表达式进行符号运算.面对繁琐、机械的计算时（例如对一个比值求二阶导数），它们尤其有用.本附录通过各种问题说明这类系统的使用.理解和使用统计学绝不以使用计算机代数系统为必要条件；不过，它不仅能减轻一些枯燥的计算工作，还能引导我们得到新的认识和更实用的答案.

需要认识到，计算机代数系统有很多种，它们可以帮助完成的计算也多种多样（求和、积分、模拟等等）.本附录的目的既不是教授这些系统的用法，也不是展示它们的一切可能用途，而只是说明其中一些可能性.

我们使用Mathematica软件包演示计算.其他软件包，例如Maple，也能用于这些计算.

\needspace{5\baselineskip}\medskip\noindent\textit{\hyperref[sec:1]{第1章}}
\begin{example}[无序抽样]\label{ex:12.6.1}
我们给出Mathematica代码，枚举\autoref{ex:1.2.20}中从$\{2,4,9,12\}$有放回抽样得到的无序结果.枚举结果并计算多项式权重后，将结果与权重排序.注意，要画出\autoref{fig:1.2.2}中的直方图，还需要做一些工作.例如，有两个不同结果的平均值都是8，因此要得到\autoref{fig:1.2.2}那样的图形，需要将结果$\{8,\frac{3}{128}\}$、$\{8,\frac{3}{64}\}$合并成$\{8,\frac{9}{128}\}$.

集合中的数超过七个时，这种枚举会变得非常耗时，因为这会产生$\binom{13}{7}=27,132$个无序结果.\footnote{译注：原文等式照录，但$\binom{13}{7}=1716$，与所列$27,132$不一致.}
\begin{enumerate}[label=(\arabic*)]
\item “DiscreteMath”包包含用于计算排列与组合的函数.
\begin{sourcecode}
In[1]:= Needs["DiscreteMath`Combinatorica`"]
\end{sourcecode}
\item 令$x=$数值的集合.不同样本的个数为$\operatorname{NumberofCompositions}[n,m]=\binom{n+m-1}{n}$.
\begin{sourcecode}
In[2]:= x ={2,4,9,12};
n=Length[x];
ncomp=NumberOfCompositions[n,n]
\end{sourcecode}
\noindent$\textbf{Out[4]}=35$.
\item 枚举样本（w），计算每个样本的平均值（avg），并计算各个值的权重（wt）.权重就是对应于该配置的多项式系数.
\begin{sourcecode}
In[5]:= w = Compositions[n, n];
wt = n!/(Apply[Times, Factorial /@ w, 1]*n^n);
avg = w . x/n;
Sort[Transpose[{avg, wt}]]
Out[8] = {2, 1/256}, {5/2, 1/64}, {3, 3/128}, {7/2, 1/64}, {15/4, 1/64},
{4, 1/256}, {17/4, 3/64}, {9/2, 1/64}, {19/4, 3/64}, {5, 3/64},
{21/4, 1/64}, {11/2, 3/128}, {11/2, 3/64}, {6, 1/64}, {6, 3/64},
{25/4, 3/64}, {13/2, 3/128}, {27/4, 3/32}, {7, 3/128}, {29/4, 1/64},
{29/4, 3/64}, {15/2, 3/64}, {31/4, 1/64}, {8, 3/128}, {8, 3/64},
{17/2, 3/64}, {35/4, 3/64}, {9, 1/256}, {37/4, 3/64}, {19/2, 1/64},
{39/4, 1/64}, {10, 1/64}, {21/2, 3/128}, {45/4, 1/64}, {12, 1/256}
\end{sourcecode}
\end{enumerate}
\end{example}
\sourcepageanchor{529}
\needspace{5\baselineskip}\medskip\noindent\textit{\hyperref[sec:2]{第2章}}
\begin{example}[一元变换]\label{ex:12.6.2}
\hyperref[exercise:2.1]{习题2.1(a)}是一个标准的一元变量变换.对计算机代数程序来说，这类计算通常很容易.
\begin{enumerate}[label=(\arabic*)]
\item 输入$f(x)$，并解出变换后的变量.\footnote{译注：代码中的反斜线以及最后一个复根的符号、$y$的幂次均按原页保留.最后一项一般不满足$x^3=y$；本例后续仅使用第一个实根$x=y^{1/3}$.本译稿没有执行或修补这些历史代码.}
\begin{sourcecode}
In[1]:=f[x_] := 42*(x^5)*(1 - x)
sol = Solve[y == x\^3, x]
\end{sourcecode}
\begin{equation*}
\textbf{Out[2]}=\{\{x\mathbin{->}y^{1/3}\},\{x\mathbin{->}-(-1)^{1/3}y^{1/3}\},\{x\mathbin{->}-(-1)^{2/3}y^{2/3}\}\}.
\end{equation*}

\item 计算雅可比因子.
\begin{sourcecode}
In[3]:= D[x/.sol[[1]],y]
\end{sourcecode}
$$\textbf{Out[3]}=\frac{1}{3y^{2/3}}.$$
\item 计算变换后变量的密度.
\begin{sourcecode}
In[4]:= f[x/.sol[[1]]]* D[x/.sol[[1]],y]
\end{sourcecode}
$$\textbf{Out[4]}=14(1-y^{1/3})y.$$
\end{enumerate}
\end{example}

\needspace{5\baselineskip}\medskip\noindent\textit{\hyperref[sec:4]{第4章}}
\begin{example}[二元变换]\label{ex:12.6.3}
下面演示一些二元变换.使用类似的代码，也可以完成多元变换.第一个计算演示\autoref{ex:4.3.4}，求正态变量之和的分布.然后计算\autoref{ex:4.3.3}：先对两个贝塔密度的乘积进行二元变换，再求边缘分布.
\begin{enumerate}[label=(\arabic*)]
\item 正态变量之和.Out[4]是联合分布，Out[5]是边缘分布.\footnote{译注：原代码对$v$从0积到$\infty$，只覆盖其支持的一半，因此所列Out[5]是完整边缘密度的一半.若求$u=X+Y$的密度，需对$v$从$-\infty$积到$\infty$.}
\begin{sourcecode}
In[1]:=f[x_,y_]:=(1/(2*Pi))*E^(-x^2/2)*E^(-y^2/2)
So:=Solve[{u==x+y,v==x-y},{x,y}]
g:=f[x/.So,y/.So]*Abs[Det[Outer[D,First[{x,y}/.So],{u,v}]]]
Simplify[g]
\end{sourcecode}
$$\textbf{Out[4]}=\left\{\frac{e^{-u^2/4-v^2/4}}{4\pi}\right\}]$$
\begin{sourcecode}
In[5]:= Integrate[g,{v,0,Infinity}]
\end{sourcecode}
$$\textbf{Out[5]}=\left\{\frac{e^{-u^2/4}}{4\sqrt{\pi}}\right\}]$$
\item 贝塔变量的乘积.（“ContinuousDistributions”包包含许多标准分布的pdf和cdf.）Out[10]是贝塔变量乘积的联合密度，Out[11]是$u$的密度.If语句读作If(检验条件,真时的值,假时的值)，因此检验条件为真时，取中间的值.在大多数情形下，检验条件为真，边缘密度就是所给的贝塔密度.\footnote{译注：由变换$u=xy,v=x$及$0<x,y<1$，联合支持应为$0<u<v<1$；原代码的积分写成$v\in(0,1)$，没有显式编码这一支持限制.原输出中还存在括号不配对，以下均照录，不将其当作已核验的可执行代码.}
\sourcepageanchor{530}
\begin{sourcecode}
In[6]:= Needs["Statistics`ContinuousDistributions`"]
Clear[f, g, u, v, x, y, a, b, c]
f[x_, y_] := PDF[BetaDistribution[a, b],x]
*PDF[BetaDistribution[a + b, c], y]
So := Solve[{u == x*y, v == x}, {x, y}]
g[u_, v_] = f[x /. So, y /. So]
*Abs[Det[Outer[D, First[{x, y} /. So], {u, v}]]]
Integrate[g[u, v] ,{v,0,1}]
\end{sourcecode}
$$\textbf{Out[10]}=\left\{\frac{(1-u/v)^{-1+c}(1-v)^{-1+b}(u/v)^{-1+a+b}v^{-1+a}}{\operatorname{Abs}[v]\operatorname{Beta}[a,b]\operatorname{Beta}[a+b,c]}\right\}$$
\begin{equation*}\begin{aligned}
\textbf{Out[11]}={}&\{\operatorname{If}[\operatorname{Re}[b]>0\mathbin{\&\&}\operatorname{Re}[b+c]<1\mathbin{\&\&}\operatorname{Im}[u]==0\mathbin{\&\&}u>0\mathbin{\&\&}u<1),\\
&\frac{1}{\operatorname{Gamma}[1-c]}\Bigl((1-u)^{-1+b+c}(-1/u)^{b+c}(-u)^cu^{-1+a+b}\\
&\hspace{3em}\operatorname{Gamma}[b]\operatorname{Gamma}[1-b-c]\Bigr),\\
&\left.\int_0^1\frac{(1-u/v)^{-1+c}(1-v)^{-1+b}(u/v)^{-1+a+b}v^{-1+a}}{\operatorname{Abs}[v]\operatorname{Beta}[a,b]\operatorname{Beta}[a+b,c]}\,dv\right\}.
\end{aligned}\end{equation*}
\end{enumerate}
\end{example}

\begin{example}[正态概率]\label{ex:12.6.4}
\hyperref[exercise:4.14]{习题4.14(a)}所要求的计算很容易处理.首先直接计算；Mathematica容易以数值方式求出积分，却没有找到闭式表达式.注意，答案用erf函数表示，该函数定义为
$$\operatorname{erf}(z)=\frac{2}{\sqrt{\pi}}\int_0^z e^{-t^2}\,dt.$$
如果用一个卡方随机变量重新表达此概率，就能找到闭式表达式.

为求积分，先设置被积函数与积分限.
\begin{sourcecode}
In[1]:= Needs["Statistics`ContinuousDistributions`"]
Clear[f, g, x, y]
f[x_, y_] = PDF[NormalDistribution[0, 1],x]
*PDF[NormalDistribution[0, 1],y]
g[x_] = Sqrt[1 - x^2]
\end{sourcecode}
$$\textbf{Out[3]}=\frac{e^{-x^2/2-y^2/2}}{2\pi}.$$
现在计算二重积分，得到Erf函数.命令N[\%]对前一行作数值求值.
\sourcepageanchor{531}
\begin{sourcecode}
In[5]:= Integrate[f[x, y], {x, -1, 1}, {y, -g[x], g[x]}]
N[%]
\end{sourcecode}
$$\textbf{Out[5]}=\frac{\int_{-1}^1e^{-x^2/2}\operatorname{Erf}\!\left(\frac{\sqrt{1-x^2}}{\sqrt2}\right)\,dx}{\sqrt{2\pi}},\qquad\textbf{Out[6]}=0.393469.$$
当然，我们知道$X^2+Y^2$是自由度为2的卡方随机变量.利用这一事实，就得到闭式答案.
\begin{sourcecode}
In[7]: = Clear[f, t]
f[t_] = PDF[ChiSquareDistribution[2], t];
Integrate[f[t], {t, 0, 1}]
N[%]
\end{sourcecode}
$$\textbf{Out[10]}=\frac12\left(2-\frac{2}{\sqrt e}\right),\qquad\textbf{Out[11]}=0.393469.$$
\end{example}
\needspace{5\baselineskip}\medskip\noindent\textit{\hyperref[sec:5]{第5章}}
\begin{example}[和的密度]\label{ex:12.6.5}
用于说明\autoref{thm:5.2.9}的\autoref{ex:5.2.10}，其中的计算相当复杂.下面在正态、柯西和Student的$t$这三种情形中演示这种计算.

有两点需要注意.
\begin{enumerate}[label=(\arabic*)]
\item 为正确解释答案，可能需要一些复分析知识.正态情形下，报告的答案带有关于变量$z$实部取值的条件（$z$可能取复数值）.在柯西例子中，重要的是知道$I^2=-1$，因此有
$$\frac{2}{\pi(-2I+z)(2I+z)}=\frac{2}{\pi(4+z^2)}.$$
\item Student的$t$变量相加时，似乎当自由度之和为偶数时存在闭式表达式，否则积分就必须用数值方法求解.（这是我们试用计算机代数系统时发现的一条经验观察.）
\end{enumerate}
还应指出，较新版本的计算机代数程序也许会避免这里出现的复数.不过，其他计算中仍会出现复数，因此最好做好处理它们的准备.
\begin{enumerate}[label=(\arabic*)]
\item 两个正态变量之和的密度.\footnote{译注：原页中$f$的代码括号不配对，分母写为$2\pi$；所列Out[3]也没有给出通常的$N(0,2)$密度$e^{-z^2/4}/(2\sqrt\pi)$.这里保留原代码和输出，不据其报告实际运行结果.}
\begin{sourcecode}
In[1]:= Clear[f, x, y, z]
f[x_] = Exp[(-x^2/2]/(2 Pi);
Integrate[f[y]* f[z - y], {y,-Infinity,Infinity}]
\end{sourcecode}
$$\textbf{Out[3]}=\operatorname{If}\left[\operatorname{Re}[z]<0,\frac{E^{-z^2/4}}{2\pi},\int_{-\infty}^{\infty}\frac{E^{-y^2/2-\frac12(-y+z)^2}}{2\pi}\,dy\right].$$
\item 两个柯西变量之和的密度.
\begin{sourcecode}
In[4]:= Clear[f, x, y, z]
f[x_] = 1/(Pi*(1+x^2));
Integrate[f[y]* f[z - y], {y,-Infinity,Infinity}]
\end{sourcecode}
$$\textbf{Out[6]}=\frac{2}{\pi(-2I+z)(2I+z)}.$$
\sourcepageanchor{532}
\item 两个自由度为5的$t$变量之和的密度.\footnote{译注：下面原代码的PDF函数少了一个右方括号；规范的函数调用应为\texttt{PDF[StudentTDistribution[5], x]}.原排印保留.}
\begin{sourcecode}
In[7]:= Needs["Statistics`ContinuousDistributions`"]
Clear[f, x, y, z]
f[x_] = PDF[StudentTDistribution[5, x]
Integrate[f[y]* f[z - y], {y,-Infinity,Infinity}]
\end{sourcecode}
$$\textbf{Out[10]}=\frac{400\sqrt5(8400+120z^2+z^4)}{3\pi(20+z^2)^5}.$$
\end{enumerate}
\end{example}

\begin{example}[均匀变量之和的四阶矩]\label{ex:12.6.6}
\autoref{exercise:5.51}要求计算12个均匀随机变量之和的四阶矩.由于密度具有分段形式，推导它有些费力（不过也可以用计算机代数完成）.然而，使用矩母函数可以简化事情.
\begin{enumerate}[label=(\arabic*)]
\item 先计算均匀随机变量$X_1$的矩母函数，再计算$\sum_{i=1}^{12}X_i$的矩母函数，其中各$X_i$相互独立.
\begin{sourcecode}
In[1]:= M[t_] = Integrate[Exp[t*x], {x, 0, 1}]
In[2]:= Msum[t_] = M[t]^12
\end{sourcecode}
$$\textbf{Out[1]}=\frac{-1+e^t}{t},\qquad\textbf{Out[2]}=\frac{(-1+e^t)^{12}}{t^{12}}.$$
\item 计算$\sum_{i=1}^{12}X_i-6$的矩母函数的四阶导数.它太长，不适合打印出来.
\begin{sourcecode}
In[3]:= g[t_]=D[Exp[-6*t]*Msum[t],{t,4}];
\end{sourcecode}
\item 值g[0]就是四阶矩.但是直接代入0会导致除以零，因此必须用极限进行计算.
\begin{sourcecode}
In[4]:= g[0]
\end{sourcecode}
\noindent Power:infy: Infinite expression $\frac{1}{0^{16}}$ encountered.
\begin{sourcecode}
In[5]:= Limit[g[t],t->0]
\end{sourcecode}
$$\textbf{Out[5]}=\frac{29}{10}.$$
\end{enumerate}
\end{example}

\needspace{5\baselineskip}\medskip\noindent\textit{\hyperref[sec:7]{第7章}}\footnote{译注：此处原小标题为“Chapter 7”，但正文引用的是第10章的例与图，以下保留原小标题.}
\begin{example}[伽马均值的渐近相对效率]\label{ex:12.6.7}
\autoref{ex:10.1.18}中得到\autoref{fig:10.1.1}的计算是在Mathematica中完成的.下面的代码会生成其中一条渐近相对效率（ARE）曲线.
\begin{enumerate}[label=(\arabic*)]
\item 用符号方法求对数似然的二阶导数.
\begin{sourcecode}
In[1]:= Needs["Statistics`ContinuousDistributions`"]
Clear[m, b, x]
f[x_, m_, b_] = PDF[GammaDistribution[m/b, b], x];
loglike2[m_, b_, x_] = D[Log[f[x, m, b]], {m, 2}];
\end{sourcecode}
\item 对二阶导数按该密度积分，求出渐近方差.
\begin{sourcecode}
In[5]:= var[m_, b_] := 1/(-Integrate[loglike2[m, b, x]*f[x, m, b],
 {x, 0, Infinity}])
\end{sourcecode}
\sourcepageanchor{533}
\item 以下代码设置绘图.
\begin{sourcecode}
In[6]:=
mu = {1, 2, 3, 4, 6, 8, 10};
beta = 5;
mlevar = Table[var[mu[[i]], beta], {i, 1, 7}];
momvar = Table[mu[[i]]*beta, {i, 1, 7}];
ARE = momvar/mlevar
Out[10] = {5.25348, 2.91014, 2.18173, 1.83958, 1.52085, 1.37349, 1.28987}
In[11]:=ListPlot[Transpose[{mu, ARE}], PlotJoined -> True,
 PlotRange -> {{0, mu[[7]]}, {0, 8}}, AxesLabel -> {"Gamma mean",
 "ARE"}]
\end{sourcecode}
\end{enumerate}
\end{example}

\needspace{5\baselineskip}\medskip\noindent\textit{\hyperref[sec:9]{第9章}}
\begin{example}[卡方矩母函数的极限]\label{ex:12.6.8}
计算\hyperref[exercise:9.30]{习题9.30(b)}中的极限分布需要小心，但在Mathematica中相当直接.
\begin{enumerate}[label=(\arabic*)]
\item 先计算自由度为$n$的卡方随机变量的矩母函数.（当然，这一步其实不必要.）\footnote{译注：此处Out[3]中外部的$2^{n/2}$因子和中间项的正指数均与卡方矩母函数不符，原页照录.满足条件时应为$(1-2t)^{-n/2}$；下面作者定义M[t\_]时使用了这个正确形式.}
\begin{sourcecode}
In[1]:= Needs["Statistics`ContinuousDistributions`"]
f[x_] = PDF[ChiSquareDistribution[n], x];
Integrate[Exp[t*x]*f[x], {x, 0, Infinity}]
\end{sourcecode}
\begin{equation*}\begin{aligned}
\textbf{Out[3]}={}&\frac{2^{n/2}}{\operatorname{Gamma}[n/2]}\operatorname{If}\Biggl[\operatorname{Re}[n]>0\mathbin{\&\&}\operatorname{Re}[t]<\frac12,\\
&\left(\frac12-t\right)^{n/2}\operatorname{Gamma}[n/2],\int_0^\infty e^{-x/2+tx}x^{-1+n/2}\,dx\Biggr].
\end{aligned}\end{equation*}

\item 因为检验条件得到满足，$\chi_n^2$的矩母函数就是中间项.现在求$\frac{\chi_n^2-n}{\sqrt{2n}}$的矩母函数的极限.
\begin{sourcecode}
In[4]:M[t_] = (1 - 2*t)^(-n/2);
Limit[Exp[-n*t/Sqrt[2*n]]*M[t/Sqrt[2*n]], n -> Infinity]
\end{sourcecode}
$$\textbf{Out[5]}=e^{t^2/2}.$$
\end{enumerate}
\end{example}

\endgroup
